\chapter{System Design and Methodology}

\section{Training Loop}
\ \ \ \ \, The simulator follows the standard synchronous FL round structure.
At round $t$, the server broadcasts the global model $w^t$ to all
clients. Each client copies the global model, trains locally for $E$
epochs using SGD with local learning rate $\eta_{\rm loc}$, and returns
a pseudo-gradient
\begin{equation}
    \Delta_n^t = w^t - [w_n^t]^{(E)},
\end{equation}
where $[w_n^t]^{(E)}$ is the local weight after $E$ epochs. Because the
local update reduces $F_n$, we have
$\Delta_n^t \approx \eta_{\rm loc}\sum_{e=1}^{E} \nabla F_n(w_{n,e-1}^t)$;
that is, $\Delta_n^t$ points in the same direction as the accumulated
local gradient, so the server update
$w_{t+1} = w_t - \eta\, \tilde{g}^t$ moves along a descent
direction of $F$ up to channel noise. The OTA oracle aggregates these
pseudo-gradients:
\begin{equation}
    \tilde{g}^t =
    \frac{1}{N\eta_{\rm loc}}
    \left(\sum_{n=1}^{N} \Delta_n^t + \xi^t\right),
    \quad
    \xi^t \sim S\alpha S(\gamma),
\end{equation}
where $S\alpha S(\gamma)$ denotes a symmetric $\alpha$-stable random
vector with tail index $\alpha$ and scale $\gamma$. The division by
$\eta_{\rm loc}$ converts the local parameter difference back into an
average-gradient scale, matching the implementation in
\path{src/experiments/federated.py}. The special case $\alpha=2$
corresponds to AWGN with standard deviation $\sqrt{2}\,\gamma$ under the
sampler convention. The server optimizer receives $\tilde{g}^t$ and
updates the global model in place.

\section{Median Anchored Clipping}
The simulator includes Median Anchored Clipping (MAC) as a robust
pre-processing step before the server optimizer. For an aggregated vector
$\tilde{g}^t\in\mathbb{R}^d$, MAC anchors the update at the coordinate-wise
median $m$, sets the clipping radius to $k$ times the Median Absolute
Deviation (MAD), and projects every coordinate into the resulting band:
\begin{align}
    m   &= \mathrm{median}\bigl\{\tilde{g}_i^t : i\in\{1,2,\dots,d\}\bigr\}, \\
    \tau &= k\cdot \mathrm{median}\bigl\{\,\lvert\tilde{g}_i^t-m\rvert : i\in\{1,2,\dots,d\}\bigr\}, \\
    \tilde{g}_i^t &\leftarrow m + \mathrm{clip}\bigl(\tilde{g}_i^t-m,\,-\tau,\,+\tau\bigr).
\end{align}
The first equation defines the anchor, the second sets the clipping radius
$\tau$ as $k$ times the MAD of the aggregate, and the third clips each
coordinate's deviation from the anchor to the band $[-\tau,\tau]$ and adds
the anchor back. $k$ here serves as the clipping factor with default setting $k=3$,
which matches the implementation in \path{src/channel/mac.py}. MAC is
expected to have little effect in the Gaussian endpoint, but it can reduce
rare impulsive coordinates when $\alpha<2$.

\section{Server-Side Baselines}
The plain OTA baseline applies the noisy aggregate directly:
\begin{equation}
    w_{t+1} = w_t - \eta \tilde{g}^t.
\end{equation}

The momentum baseline maintains a first-moment buffer:
\begin{equation}
    m_t = \rho m_{t-1} + \tilde{g}^t,
    \quad
    w_{t+1} = w_t - \eta m_t.
\end{equation}

These baselines are useful because they isolate whether the second-moment
normalization in adaptive methods provides additional robustness beyond
simple smoothing.

\section{AdaGrad-OTA}
AdaGrad-OTA maintains a smoothed update direction $s_t$ and an accumulated
coordinate-wise magnitude $v_t$:
\begin{align}
    s_t &= \beta_1 s_{t-1} + (1-\beta_1)\tilde{g}^t, \\
    v_t &= v_{t-1} + |s_t|^{\alpha}, \\
    w_{t+1} &= w_t -
    \eta \frac{s_t}{v_t^{1/\alpha}+\varepsilon}.
\end{align}
For $\alpha=2$, the denominator reduces to the usual square-root AdaGrad
normalization. For $\alpha<2$, the same form follows the fractional
$\alpha$-stable accumulation rule used in the reference paper. Because $v_t$ is
monotone increasing, AdaGrad-OTA tends to become conservative over long runs.
This can stabilize noisy training, but it may also slow late-stage
convergence.

\section{Adam-OTA}
The implemented Adam-OTA variant uses a stable server-side Adam update on the
OTA-aggregated direction:
\begin{align}
    m_t &= \beta_1 m_{t-1} + (1-\beta_1)\tilde{g}^t, \\
    v_t &= \beta_2 v_{t-1} + (1-\beta_2)|\tilde{g}^t|^{\alpha}, \\
    \hat{m}_t &= \frac{m_t}{1-\beta_1^t}, \\
    \hat{v}_t &= \frac{v_t}{1-\beta_2^t}, \\
    w_{t+1} &= w_t -
    \eta \frac{\hat{m}_t}{\hat{v}_t^{1/\alpha}+\varepsilon}.
\end{align}
When $\alpha=2$, this is classical Adam applied at the server to the
noisy OTA aggregate. For $\alpha<2$, the second-moment update becomes a
fractional moment accumulator. This preserves the project hypothesis
because the update uses a state variable to compress effective step
sizes when the observed aggregate has large magnitude.

\paragraph{Implementation note: $\varepsilon$.}
In the code, AdaGrad-OTA uses $\varepsilon=10^{-4}$ while Adam-OTA uses
$\varepsilon=10^{-8}$. The two differ by four orders of magnitude by
design: AdaGrad-OTA has no $\beta_2$ exponential moving average, so its
first-round $v_t$ is very small and a larger $\varepsilon$ is needed to
damp the early transient; Adam-OTA's bias correction
$\hat v_t = v_t/(1-\beta_2^{\,t})$ delivers a non-negligible denominator
already after a few rounds, so the smaller textbook default
$\varepsilon=10^{-8}$ is kept. The asymmetry between the two optimizers
is the dominant reason why their early-round test-loss transients in
Figure~\ref{fig:cifar-curves} differ in shape.

\section{Finite-Trace Stability Analysis}
\label{sec:finite-trace-stability}
The project does not reproduce the full non-convex convergence proof in
\cite{wang2024adaptive}. It does, however, admit a rigorous finite-trace
analysis that matches the implemented simulator. The analysis is
coordinate-wise and concerns the denominator that turns a noisy OTA
aggregate into an adaptive server step. For a fixed coordinate $i$, let
$z_t$ denote the quantity accumulated by the optimizer: $z_t=s_{t,i}$ for
AdaGrad-OTA and $z_t=\tilde g^t_i$ for Adam-OTA. Define $x_t=|z_t|$.

\paragraph{Moment caveat.}
For a non-Gaussian symmetric $\alpha$-stable random variable
($0<\alpha<2$), $\mathbb{E}|\xi|^p$ is finite only for $0<p<\alpha$ and
is infinite for $p\geq\alpha$~\cite{nikias1995alphastable}; the endpoint
$\alpha=2$ is the Gaussian case. Therefore, in the heavy-tailed regime
$\alpha<2$, the usual bounded-variance SGD argument is not the right tool,
and even the population $\alpha$-moment used by the fractional accumulator
need not exist. The implemented algorithms instead operate on a finite
realized training trace, where every stored accumulator is finite unless
numerical overflow occurs. The propositions below should be read in that
finite-trace sense.

\paragraph{Proposition 1: impulse contraction.}
For AdaGrad-OTA, define
\begin{equation}
    V_t=\sum_{\tau=1}^{t} x_\tau^\alpha,
    \qquad
    \eta_{t,i}^{\rm eff}=\frac{\eta}{V_t^{1/\alpha}+\varepsilon}.
\end{equation}
If round $t$ contains an impulsive coordinate with magnitude $x_t=R$, then
\begin{equation}
    \eta_{t,i}^{\rm eff}\leq \frac{\eta}{R+\varepsilon},
    \qquad
    \left|\eta\frac{z_t}{V_t^{1/\alpha}+\varepsilon}\right|\leq \eta.
\end{equation}
Thus the same observation that creates a large numerator also enlarges the
denominator enough to cap the coordinate-wise normalized step. A fixed-step
FedAvg-OTA update has no such feedback: its coordinate response to the same
impulse has magnitude $\eta R$.

\noindent\emph{Proof.}
Since $V_t=V_{t-1}+R^\alpha\geq R^\alpha$, we have
$V_t^{1/\alpha}\geq R$. Substituting this inequality into the effective
learning rate gives the first bound, and multiplying by $R$ gives the
normalized-step bound.

For Adam-OTA, the bias-corrected second accumulator can be written as
\begin{equation}
    \hat V_t
    = \sum_{\tau=1}^{t}
    \frac{(1-\beta_2)\beta_2^{t-\tau}}{1-\beta_2^t}
    x_\tau^\alpha .
\end{equation}
The current sample has weight
$(1-\beta_2)/(1-\beta_2^t)\geq 1-\beta_2$. Therefore, after an impulse
$x_t=R$,
\begin{equation}
    \frac{\eta}{\hat V_t^{1/\alpha}+\varepsilon}
    \leq
    \frac{\eta}{(1-\beta_2)^{1/\alpha}R+\varepsilon}.
\end{equation}
Adam-OTA keeps a finite memory rather than a monotone memory, so this bound
contains the constant $(1-\beta_2)^{1/\alpha}$, but the same denominator
feedback remains present.

\paragraph{Proposition 2: tail-index step-size scaling.}
For AdaGrad-OTA, suppose that along a finite trace the empirical
$\alpha$-power average is bounded below by a positive constant,
\begin{equation}
    c_t=\frac{1}{t}\sum_{\tau=1}^{t}x_\tau^\alpha \geq c>0,
    \qquad 1\leq t\leq T.
\end{equation}
Then each coordinate-wise effective learning rate satisfies
\begin{equation}
    \eta_{t,i}^{\rm eff}
    \leq
    \eta c^{-1/\alpha}t^{-1/\alpha},
\end{equation}
and its average over the first $T$ rounds obeys
\begin{equation}
    \frac{1}{T}\sum_{t=1}^{T}\eta_{t,i}^{\rm eff}
    \leq
    \eta c^{-1/\alpha}
    \left(\frac{1}{T}+\frac{\alpha}{\alpha-1}T^{-1/\alpha}\right).
\end{equation}

\noindent\emph{Proof.}
The denominator satisfies
$V_t^{1/\alpha}=(t c_t)^{1/\alpha}\geq (tc)^{1/\alpha}$, so
$\eta_{t,i}^{\rm eff}\leq\eta c^{-1/\alpha}t^{-1/\alpha}$. Averaging this
bound and using
$\sum_{t=1}^{T}t^{-1/\alpha}\leq
1+\int_1^T u^{-1/\alpha}\,du
\leq 1+\frac{\alpha}{\alpha-1}T^{1-1/\alpha}$
gives the result.

This proposition gives a direct tail-index interpretation. Smaller
$\alpha$ produces the faster decay exponent $1/\alpha$, so heavy-tailed
traces force stronger automatic damping. That damping improves stability
after impulsive rounds, but it can also compress useful signal and slow
late-stage optimization.

\paragraph{One-dimensional interpretation.}
The same mechanism can be visualized with the quadratic problem
\begin{equation}
    f(w)=\frac{1}{2}(w-w^\star)^2.
\end{equation}
The noiseless gradient is $w_t-w^\star$, while the OTA server observes
\begin{equation}
    \tilde{g}_t = w_t-w^\star+\xi_t,
    \quad \xi_t \sim S\alpha S(\gamma).
\end{equation}
For an AdaGrad-style update,
\begin{equation}
    v_t = v_{t-1}+|\tilde{g}_t|^\alpha,
    \quad
    w_{t+1}=w_t-\eta\frac{\tilde{g}_t}{v_t^{1/\alpha}+\epsilon}.
\end{equation}
Whenever an impulsive noise sample makes $|\tilde{g}_t|$ large, it also
makes $v_t$ larger. The next effective step size
\begin{equation}
    \eta_t^{\mathrm{eff}} =
    \frac{\eta}{v_t^{1/\alpha}+\epsilon}
\end{equation}
therefore contracts automatically. This special case is not a replacement
for the full theorem in \cite{wang2024adaptive}, but it explains why the
implemented accumulator reacts strongly after impulsive rounds.

\section{Implementation Modules}
The codebase is organized around separable modules. The channel layer
contains the noisy aggregation oracle and noise sampler. The data layer
contains dataset loading and client partitioning. The model layer contains
classification networks. The optimizer layer contains the four server
optimizers. The experiment layer contains single-run training, method
comparison, ablation, plotting, logging and result export.

This organization supports the project goal in two ways. First, the
$\alpha$-stable experiments and MAC comparisons can be run reproducibly through
command-line scripts. Second, the same interfaces can later be used to add
richer physical channel models, additional datasets, or stricter reproduction
settings without rewriting the training loop.
